\section{Conclusões}
%%%%%%%%%%%%%%%%%%%%%%%%%%
\begin{frame}[fragile]
\frametitle{Conclusões}
\begin{minipage}{0.47\textwidth}
    \begin{itemize}
        \item O que foi visto
        \item O que tem a ser feito
        \item Oportunidades
    \end{itemize}
\end{minipage}
\begin{minipage}{0.5\textwidth}
\begin{figure}[ht!]
\begin{center}
\includegraphics[width=1.2\textwidth, height=0.40\textheight]{figures/logo_picat_alex.jpg}
\end{center}
\end{figure}
\end{minipage}
\end{frame}
%%%%%%%%%%%%%%%%%%%%%%%%%%
\begin{frame}[fragile]

    \frametitle{Conclusões}

    \begin{itemize}
      \item Picat é jovem  $\approx$ 2013
      \pause
      \item Uma evolução ao Prolog após seus mais de 40 anos de existência e sucesso!
      \pause
      \item Sua sintaxe é moderna e intuitiva
      \pause
      \item Código aberto, multi-plataforma e repleta de possibilidades (sistemas embarcados, combinatória,
      SO, escalonamentos, planejamento etc) $\Rightarrow$ NP's completos em geral
%      \pause
%      \item Uso para fins diversos
      \pause
      \item Muitas bibliotecas específicas prontas: CP (PR), SAT, Planner, etc
      \pause
      \item A sintaxe de PR exige um pouco mais do programador
      \pause
      \item Dúvidas: o guia do usuário, livro do Hakan e o Fórum de discussão do Picat

    \end{itemize}
\end{frame}


%%%%%%%%%%%%%%%%%%%%%%%%%%

\subsection{Ficou Faltando}
\begin{frame}[fragile]

    \frametitle{Ficou Faltando:}

    \begin{itemize}
      \item Uso do \texttt{debug} e \texttt{trace} (cansativo -- uma oportunidade) $\Rightarrow$ \textcolor{magenta}{\textbf{Pendência resolvida!}} 

     \pause
     \item Explorar uso dos \textit{solvers} de PO (fácil)

      \pause
      \item Explorar a criação e  uso de módulos  (mais fácil ainda)



    \end{itemize}
\end{frame}



\subsection{Dicas de Programação}
\begin{frame}[fragile] 

    \frametitle{Dicas de Programação:}

    \begin{itemize}
      \item Use o interpretador e o compilador concomitamente. O interpretador
      sempre acusa \textit{warnings} etc. O modo compilado na console não apresenta
      alguns \textit{warnings}. Efeito \textit{cascata} de erro ...

      \pause
      \item No modo interpretado, cada linha de código pode ser testada isoladamente,
      assim, o efeito global desta é restrita. Qualquer erro ou falha é rapidamente
      detectada.

       \pause
      \item Consulte o manual do usuário \textit{on-line} em \textit{html}, 
      mantido pelo Alexandre Gonçalves -- UFSC
      \url{http://retina.inf.ufsc.br/picat_guide}


      \pause
      \item Consulte o site do Picat e dos grandes mestres Hakan, Neng-Fa, Roman Barták, 
      Sergii Dymchenko, etc

      \pause
      \item Inscreva-se no fórum e consulte o Guia do Usuário (tudo em inglês)


    \end{itemize}
\end{frame}


\subsection{Agradecimentos}
\begin{frame}[fragile]

    \frametitle{Agradecimentos}

    \begin{itemize}
      \item Muito obrigado a voces!
      \pause
%      \item Algumas pessoas que deram opiniões neste material 
%            e me incentivaram a terminá-lo 
%      \pause%
%      \item  Em especial a minha esposa: Maria Luiza
        \item  Ao Glauco pelo convite
        
      \pause
      \item Claudio Cesar de Sá
     \item Contacto: 
     \begin{description}
     	\item[\Letter \/] \url{ccs1664@gmail.com}
     	\item[\Letter \/] \url{ccs1664@yahoo.com}
         	\item \url{https://claudiocesar.wordpress.com/}     	
        %  \item[\Letter \/] \url{claudio.sa@udesc.br}
         % \item[\Letter \/] \url{claudio@colmeia.udesc.br}
     \end{description}


    \end{itemize}
\end{frame}



\begin{frame}[fragile]
  \frametitle{Finalizando ...}


\begin{minipage}{0.47\textwidth}

\begin{Large}
\textbf{\textcolor{magenta}{... espero que\\ tenham gostado\\ e obrigado!}}
\end{Large}

\end{minipage}
\begin{minipage}{0.5\textwidth}
\begin{figure}[ht!]
\begin{center}
\includegraphics[width=1.2\textwidth, height=0.40\textheight]{figures/logo_picat_alex.jpg}
\end{center}
\end{figure}
\end{minipage}

\end{frame}
					
